\documentclass[11pt,a4paper]{article}

% ── Packages ──────────────────────────────────────────────────────────────
\usepackage[utf8]{inputenc}
\usepackage[T1]{fontenc}
\usepackage{lmodern}
\usepackage[margin=1in]{geometry}
\usepackage{hyperref}
\usepackage{xcolor}
\usepackage{listings}
\usepackage{booktabs}
\usepackage{longtable}
\usepackage{array}
\usepackage{enumitem}
\usepackage{amsmath,amssymb}
\usepackage{parskip}
\usepackage{fancyhdr}
\usepackage{titlesec}
\usepackage{caption}
\usepackage{float}
\usepackage{graphicx}

% ── Colors ────────────────────────────────────────────────────────────────
\definecolor{codebackground}{HTML}{F5F5F5}
\definecolor{codekeyword}{HTML}{005CC5}
\definecolor{codecomment}{HTML}{6A737D}
\definecolor{codestring}{HTML}{032F62}
\definecolor{purpleballoon}{HTML}{A855F7}
\definecolor{tealballoon}{HTML}{14B8A6}
\definecolor{orangeballoon}{HTML}{F97316}

% ── Hyperref setup ────────────────────────────────────────────────────────
\hypersetup{
    colorlinks=true,
    linkcolor=blue!60!black,
    urlcolor=blue!60!black,
    citecolor=blue!60!black,
    bookmarksdepth=3
}

% ── Listings style ────────────────────────────────────────────────────────
\lstdefinestyle{pythonstyle}{
    language=Python,
    backgroundcolor=\color{codebackground},
    basicstyle=\ttfamily\small,
    keywordstyle=\color{codekeyword}\bfseries,
    commentstyle=\color{codecomment}\itshape,
    stringstyle=\color{codestring},
    breaklines=true,
    frame=single,
    rulecolor=\color{black!20},
    xleftmargin=1em,
    xrightmargin=1em,
    aboveskip=1em,
    belowskip=1em,
    showstringspaces=false,
    columns=flexible,
    morekeywords={np,stats,float,int,dict,list,tuple,bool,defaultdict}
}

\lstdefinestyle{tsstyle}{
    language=Java,
    backgroundcolor=\color{codebackground},
    basicstyle=\ttfamily\small,
    keywordstyle=\color{codekeyword}\bfseries,
    commentstyle=\color{codecomment}\itshape,
    stringstyle=\color{codestring},
    breaklines=true,
    frame=single,
    rulecolor=\color{black!20},
    xleftmargin=1em,
    xrightmargin=1em,
    aboveskip=1em,
    belowskip=1em,
    showstringspaces=false,
    columns=flexible,
    morekeywords={const,let,Math,if,return}
}

\lstdefinestyle{jsonstyle}{
    backgroundcolor=\color{codebackground},
    basicstyle=\ttfamily\scriptsize,
    breaklines=true,
    frame=single,
    rulecolor=\color{black!20},
    xleftmargin=1em,
    xrightmargin=1em,
    aboveskip=1em,
    belowskip=1em,
    showstringspaces=false,
    columns=flexible
}

% ── Header / Footer ──────────────────────────────────────────────────────
\pagestyle{fancy}
\fancyhf{}
\fancyhead[L]{\small BART Metrics Reference --- Full Technical Documentation}
\fancyhead[R]{\small\thepage}
\renewcommand{\headrulewidth}{0.4pt}

% ── Title ─────────────────────────────────────────────────────────────────
\title{\textbf{BART Metrics Reference}\\[0.3em]
\large Full Technical Documentation}
\author{Ahmet Selim Yılmaz}
\date{}

% ══════════════════════════════════════════════════════════════════════════
\begin{document}
\maketitle
\thispagestyle{fancy}

\tableofcontents
\newpage

% ──────────────────────────────────────────────────────────────────────────
\section{Explosion Model}
% ──────────────────────────────────────────────────────────────────────────

Unlike the original Lejuez et al.\ (2002) task --- which pre-draws a single breakpoint
uniformly at random and pops the balloon deterministically when it is reached --- our
variant uses a \textbf{sequential independent Bernoulli model}.
At each pump attempt $k$ the balloon explodes with probability

\begin{equation}
\label{eq:explosion}
P(\text{explode at pump } k) = \frac{k}{\texttt{maxPumps}}\,.
\end{equation}

In the client this looks like:

\begin{lstlisting}[style=tsstyle,caption={Explosion check (TypeScript, client-side).}]
const newPumps = currentBalloon.pumps + 1;
const explosionProbability = newPumps / maxPumps;
const explode = newPumps >= maxPumps
             || Math.random() < explosionProbability;
\end{lstlisting}

Two consequences are worth highlighting.
First, there is no predetermined breakpoint --- every pump is a genuinely stochastic
event, and the balloon always pops at \texttt{maxPumps}.
Second, the EV-optimal stopping point sits well below $\texttt{maxPumps}/2$:

\begin{table}[H]
\centering
\begin{tabular}{@{}llccc@{}}
\toprule
\textbf{Color} & \textbf{Hex} & \textbf{\texttt{maxPumps}} & \textbf{Risk Level} & \textbf{EV-Optimal Stop} \\
\midrule
\textcolor{purpleballoon}{Purple} & \texttt{\#A855F7} & 128 & Low    & ${\sim}\,32$ pumps \\
\textcolor{tealballoon}{Teal}     & \texttt{\#14B8A6} & 32  & Medium & ${\sim}\,8$ pumps  \\
\textcolor{orangeballoon}{Orange} & \texttt{\#F97316} & 8   & High   & ${\sim}\,2$ pumps  \\
\bottomrule
\end{tabular}
\caption{Balloon color profiles. Neutral colors were chosen deliberately to
avoid triggering common psychological associations (e.g.\ red\,$=$\,danger).}
\end{table}

\subsection{Session Generation}

Every session presents exactly 30 balloons --- 10 of each color --- shuffled with
Fisher--Yates so the participant cannot predict the sequence:

\begin{lstlisting}[style=tsstyle,caption={Shuffle logic (TypeScript).}]
["ORANGE", "TEAL", "PURPLE"].forEach((type) => {
    const profile = RISK_PROFILES[type];
    for (let i = 0; i < 10; i++)
        configs.push({ id: 0, color: profile.color,
                       maxPumps: profile.maxPumps });
});
for (let i = configs.length - 1; i > 0; i--) {
    const j = Math.floor(Math.random() * (i + 1));
    [configs[i], configs[j]] = [configs[j], configs[i]];
}
\end{lstlisting}


% ──────────────────────────────────────────────────────────────────────────
\section{Raw Event Data Format}
% ──────────────────────────────────────────────────────────────────────────

The client records every action as a timestamped JSON object.
Timestamps come from \texttt{performance.now()} (monotonic, sub-millisecond),
not \texttt{Date.now()}, to avoid wall-clock jumps.
Each event has three fields: \texttt{type}, \texttt{timestamp}, and \texttt{payload}.

\begin{lstlisting}[style=jsonstyle,caption={A few raw events from the sample session (\texttt{6fe2f79a}).}]
{"type":"pump",   "timestamp":36678.70, "payload":{"color":"purple","balloon_id":1}}
{"type":"pump",   "timestamp":37262.10, "payload":{"color":"purple","balloon_id":1}}
  ...  (4 more pumps)
{"type":"collect","timestamp":40844.90, "payload":{"color":"purple","balloon_id":1}}
{"type":"pump",   "timestamp":42964.50, "payload":{"color":"teal",  "balloon_id":2}}
{"type":"pump",   "timestamp":43418.00, "payload":{"color":"teal",  "balloon_id":2}}
{"type":"pump",   "timestamp":43836.20, "payload":{"color":"teal",  "balloon_id":2}}
{"type":"explode","timestamp":43836.30, "payload":{"color":"teal",  "balloon_id":2,
                                                    "pump_count":3}}
\end{lstlisting}

The scoring engine starts by splitting this flat list into per-balloon groups ---
a new group begins after every \texttt{collect} or \texttt{explode}:

\begin{lstlisting}[style=pythonstyle,caption={Balloon segmentation (Python).}]
def _segment_balloons(events):
    balloons, current = [], []
    for event in events:
        current.append(event)
        if event.type in ("collect", "explode"):
            balloons.append(current)
            current = []
    if current:
        balloons.append(current)
    return balloons
\end{lstlisting}

For the sample session this yields 30 balloon segments: 10~purple, 10~teal,
10~orange.


% ──────────────────────────────────────────────────────────────────────────
\section{Handling Right-Censored Data}
% ──────────────────────────────────────────────────────────────────────────

When a balloon explodes the pump count is cut short --- the participant may well
have intended to keep pumping.  This is a textbook right-censoring problem, and
ignoring it biases every mean downward.

\textbf{A concrete example.} In the sample session all 10 orange balloons exploded.
Their observed counts ($1, 4, 5, 4, 4, 4, 2, 2, 1$) average $2.9$, but
those numbers reflect where the RNG happened to terminate the trial, not the
participant's chosen stopping rule.  Without a single collected orange balloon
we simply cannot recover the true intention.

The fix is straightforward: wherever we need a measure of \emph{what the participant
chose to do}, we restrict the calculation to \textbf{collected (non-exploded)
balloons}.  A small helper manages this, falling back to all balloons (with a
warning) when fewer than two collected trials are available for a given color:

\begin{lstlisting}[style=pythonstyle,caption={Collected-preference helper (Python).}]
MIN_COLLECTED_FALLBACK = 2

def _prefer_collected(collected, all_data, min_count=2):
    if len(collected) >= min_count:
        return collected, False      # use collected
    return all_data, True            # fallback + warning
\end{lstlisting}

For orange this fallback fires routinely --- $P(\text{survive 4 pumps on orange})
\approx 16\%$, so most sessions have very few (sometimes zero) collected orange
balloons.


% ──────────────────────────────────────────────────────────────────────────
\section{Overall Metrics}
% ──────────────────────────────────────────────────────────────────────────

\subsection{\texttt{average\_pumps\_adjusted}}

The classic Lejuez ``adjusted'' BART score: mean pump count across collected
balloons only.

\begin{lstlisting}[style=pythonstyle]
non_exploded_pumps = []
for balloon_events in balloons:
    pumps = sum(1 for e in balloon_events if e.type == "pump")
    terminal = next(
        (e.type for e in reversed(balloon_events)
         if e.type in ("collect", "explode")), None)
    if terminal != "explode":
        non_exploded_pumps.append(pumps)

average_pumps_adjusted = (np.mean(non_exploded_pumps)
                          if non_exploded_pumps
                          else np.mean(all_pumps))
\end{lstlisting}

Range: $[0,\, 128]$.  In the sample session 10~balloons were collected with counts
$\{6, 5, 5, 11, 7, 3, 21, 5, 3\}$, giving $\texttt{average\_pumps\_adjusted} = 7.2$.

One important caveat: this number is colour-blind.  A participant who pumps
$5\times$ on every colour looks the same as one who pumps $20\times$ on purple
and $1\times$ on orange.  For colour-aware alternatives see
\texttt{rng\_normalized\_pumps} (\S\ref{sec:rng_norm}) or
\texttt{risk\_adjustment\_score} (\S\ref{sec:risk_adj}).


\subsection{\texttt{avg\_pumps\_all\_balloons}}

Plain mean across every balloon regardless of outcome:

\begin{equation}
    \texttt{avg\_pumps\_all\_balloons} = \frac{1}{N}\sum_{i=1}^{N} \text{pumps}_i
\end{equation}

Sample session: $168 / 30 = 5.6$.  Useful for comparisons with other BART studies
that report the unadjusted mean, but it understates behavioural intention because
exploded balloons have truncated counts.


\subsection{\texttt{explosion\_rate}}

Simply the fraction of balloons that popped:

\begin{equation}
    \texttt{explosion\_rate} = \frac{\texttt{total\_explosions}}{\texttt{total\_balloons}}
\end{equation}

Sample session: $20/30 = 0.667$.  Worth remembering that two people playing exactly
the same strategy can end up with noticeably different explosion rates on any given
run due to the stochastic nature of the task.


\subsection{\texttt{mean\_latency\_between\_pumps}}

Average milliseconds between consecutive pumps \emph{within} the same balloon.
Cross-balloon gaps (e.g.\ time spent deciding to start a new balloon) are excluded,
and any intra-balloon pause longer than 2\,s is filtered as an outlier.

\begin{lstlisting}[style=pythonstyle]
all_intra_latencies = []
for balloon_events in balloons:
    pump_times = [e.timestamp for e in balloon_events
                  if e.type == "pump"]
    if len(pump_times) >= 2:
        diffs = np.diff(pump_times)
        all_intra_latencies.extend(diffs.tolist())

latencies = np.array(all_intra_latencies)
latencies = latencies[latencies < 2000.0]
mean_latency = np.mean(latencies)
\end{lstlisting}

To make this concrete, here are the inter-pump intervals for balloon~1
(purple, collected at 6~pumps):
\begin{align*}
\text{Timestamps} &= [36678.7,\; 37262.1,\; 38028.8,\; 38628.6,\; 39362.1,\; 40012.1] \\
\text{Diffs (ms)} &= [583.4,\; 766.7,\; 599.8,\; 733.5,\; 650.0]
\end{align*}
Those five intervals enter the global pool.  Across all 30 balloons the session-level
mean is $555.4$\,ms --- on the deliberate side (fast/impulsive pumpers tend to sit
around ${\sim}200$\,ms).


\subsection{Session-Level Counts}

\texttt{total\_balloons}, \texttt{total\_pumps}, \texttt{total\_explosions}, and
\texttt{total\_collections} are straightforward aggregations.  They should always
satisfy

\begin{equation}
    \texttt{total\_explosions} + \texttt{total\_collections} = \texttt{total\_balloons}
\end{equation}

Sample session: $20 + 10 = 30$\;\checkmark.


% ──────────────────────────────────────────────────────────────────────────
\section{Per-Colour Metrics}
% ──────────────────────────────────────────────────────────────────────────

For each colour we emit a \texttt{ColorMetrics} object.  Colour is read from the
event payload (defaulting to teal if absent):

\begin{lstlisting}[style=pythonstyle]
def _extract_balloon_color(balloon_events):
    for event in balloon_events:
        if hasattr(event.payload, "color") and event.payload.color:
            return event.payload.color.lower()
    return "teal"
\end{lstlisting}

\begin{table}[H]
\centering
\begin{tabular}{@{}lp{8.5cm}@{}}
\toprule
\textbf{Field} & \textbf{Description} \\
\midrule
\texttt{color}                  & \texttt{"purple"}, \texttt{"teal"}, or \texttt{"orange"} \\
\texttt{average\_pumps}         & Mean pumps across \emph{all} balloons of this colour \\
\texttt{behavioral\_avg\_pumps} & Mean pumps from collected balloons only (falls back to all if $< 2$ collected) \\
\texttt{explosion\_rate}        & Colour-specific explosion rate \\
\texttt{total\_balloons}        & Count of balloons of this colour (expected: 10) \\
\texttt{collected\_count}       & How many survived to collection \\
\texttt{risk\_profile}          & \texttt{"low"} / \texttt{"medium"} / \texttt{"high"} \\
\texttt{used\_fallback}         & True when \texttt{behavioral\_avg\_pumps} had to fall back \\
\bottomrule
\end{tabular}
\caption{Fields in the per-colour \texttt{ColorMetrics} object.}
\end{table}

\begin{table}[H]
\centering
\begin{tabular}{@{}lccccc@{}}
\toprule
\textbf{Colour} & \textbf{Balloons} & \textbf{Collected} & \textbf{Avg Pumps (all)} & \textbf{Explosion Rate} & \textbf{Fallback?} \\
\midrule
Purple & 10 & 5 & 9.9  & 0.50 & No \\
Teal   & 10 & 5 & 4.0  & 0.50 & No \\
Orange & 10 & 0 & 2.9  & 1.00 & Yes \\
\bottomrule
\end{tabular}
\caption{Colour breakdown for the sample session.  The gap between
\texttt{behavioral\_avg\_pumps} and \texttt{average\_pumps} gets wider as the
explosion rate increases --- orange being the most extreme case.}
\end{table}


% ──────────────────────────────────────────────────────────────────────────
\section{Behavioural Intention Metrics}
% ──────────────────────────────────────────────────────────────────────────

These metrics all aim to capture what the participant \emph{intended} to do, so they
operate on collected balloons wherever possible.

\subsection{\texttt{rng\_normalized\_pumps}}
\label{sec:rng_norm}

Our primary risk measure.  For every collected balloon we divide pump count by the
colour's capacity, then average:

\begin{equation}
    \texttt{rng\_normalized\_pumps} = \frac{1}{|\mathcal{C}|}\sum_{b \in \mathcal{C}}
        \frac{\text{pumps}_b}{\text{maxPumps}_{\text{colour}(b)}}
\end{equation}

\begin{lstlisting}[style=pythonstyle]
normalized = []
for color, pumps_list in color_pumps_behavioral.items():
    cap = COLOR_PROFILES[color]["max_pumps"]
    for p in pumps_list:
        normalized.append(p / cap)

rng_normalized_pumps = np.mean(normalized)
\end{lstlisting}

Range $[0, 1]$.  A value of $0.25$ means the participant used roughly a quarter
of each balloon's capacity on average.  Because the normalisation removes the scale
difference between colours ($128$ vs $32$ vs $8$), this metric is directly comparable
across risk levels --- something \texttt{average\_pumps\_adjusted} cannot offer.


\subsection{\texttt{impulsivity\_index}}

Focuses on orange balloons specifically.  We take the mean collected-orange pump
count and normalise by the orange ceiling~(8):

\begin{equation}
    \texttt{impulsivity\_index} = \text{clip}\!\left(
        \frac{\text{mean}(\text{collected\_orange\_pumps})}{8},\; 0,\; 1\right)
\end{equation}

\begin{lstlisting}[style=pythonstyle]
orange_behavioral = color_pumps_behavioral.get("orange", [])
orange_avg = np.mean(orange_behavioral) if orange_behavioral else 0.0
impulsivity_index = np.clip(orange_avg / 8.0, 0.0, 1.0)
\end{lstlisting}

Because collected orange balloons are rare ($P(\text{survive 4 pumps}) \approx 16\%$),
each one carries a lot of signal.  A score of $0.5$ means the participant was pushing
to half of orange's capacity --- well past the EV-optimal stopping point of ${\sim}2$.

In the sample session every orange balloon exploded, so the fallback kicked in:
$2.9 / 8 = 0.36$.
(The value of $1.0$ in the attached JSON comes from an older formula based on
explosion counts; the current definition is the one above.)


\subsection{\texttt{patience\_index} and \texttt{patience\_index\_normalized}}

The mirror image of impulsivity: mean pump count on collected purple balloons,
where risk is low and the optimal strategy is to pump aggressively.

\begin{align}
    \texttt{patience\_index} &= \text{mean}(\text{collected\_purple\_pumps}) \\[2pt]
    \texttt{patience\_index\_normalized} &= \text{clip}\!\left(
        \frac{\texttt{patience\_index}}{128},\; 0,\; 1\right)
\end{align}

Sample session: $9.9$ (raw), $0.077$ (normalised) --- the participant was nowhere
near exploiting the full safety margin that purple balloons offer.


\subsection{\texttt{orange\_avg\_pumps}}

Simply the raw version of \texttt{impulsivity\_index} before normalisation ($[0, 8]$).
Sometimes easier to reason about clinically: ``the participant pumped an average of
4~times on the high-risk balloons.''


% ──────────────────────────────────────────────────────────────────────────
\section{Learning \& Adaptation Metrics}
% ──────────────────────────────────────────────────────────────────────────

Both metrics in this section try to answer the same question --- \emph{did the
participant get better over time?} --- but they approach it differently.

\subsection{\texttt{learning\_rate}}

Fits a simple OLS regression of pump count on trial number, separately per colour,
using collected balloons.  The slope is weighted by~$R^2$ so that noisy fits count
for less.  Directionality is colour-dependent: pumping \emph{more} on purple over
time is adaptive, while pumping \emph{less} on orange is.

\begin{lstlisting}[style=pythonstyle,caption={Regression-based learning rate (Python).}]
for color in color_trials_all:
    collected = color_trials_collected.get(color, [])
    all_t = color_trials_all[color]
    trials = collected if len(collected) >= 2 else all_t
    if len(trials) < 2:
        continue

    xs = np.array([t[0] for t in trials])
    ys = np.array([t[1] for t in trials])
    slope, _, r_value, _, _ = stats.linregress(xs, ys)
    weighted = slope * (r_value ** 2)

    if color == "orange":
        learning_slopes.append(-weighted)
    elif color == "purple":
        learning_slopes.append(weighted)
    else:                          # teal: half-weight
        learning_slopes.append(-weighted * 0.5)

learning_rate = np.clip(np.mean(learning_slopes), -1.0, 1.0)
\end{lstlisting}

Range $[-1, 1]$.  Sample session: $0.059$.

The main weakness is sample size.  With only 10 trials per colour --- and often
fewer collected --- a single outlier can swing the regression substantially.


\subsection{\texttt{half\_split\_learning\_rate}}

A more robust alternative that sidesteps the regression entirely.  For each colour
the collected trials are sorted chronologically and split in half; we then look at
the relative change:

\begin{equation}
    \delta_\text{colour} = \frac{\bar{x}_{\text{2nd half}} - \bar{x}_{\text{1st half}}}
                                {\bar{x}_{\text{overall}}}
\end{equation}

\begin{lstlisting}[style=pythonstyle,caption={Half-split learning rate (Python).}]
for color in color_trials_all:
    trials = collected if len(collected) >= 4 else all_trials
    if len(trials) < 4:
        continue

    sorted_t = sorted(trials, key=lambda x: x[0])
    half = len(sorted_t) // 2
    m1 = np.mean([t[1] for t in sorted_t[:half]])
    m2 = np.mean([t[1] for t in sorted_t[half:]])
    m  = np.mean([t[1] for t in sorted_t])
    if m == 0:
        continue

    delta = (m2 - m1) / m
    if color == "orange":
        scores.append(-delta)
    elif color == "purple":
        scores.append(delta)
    else:
        scores.append(-delta * 0.5)

half_split_lr = np.clip(np.mean(scores), -1.0, 1.0)
\end{lstlisting}

A value of $0.2$ means behaviour improved by roughly 20\% across the session.
Because it compares group means rather than fitting a line through a handful of
points, it is much less sensitive to individual outlier trials.


% ──────────────────────────────────────────────────────────────────────────
\section{Risk Calibration Metrics}
% ──────────────────────────────────────────────────────────────────────────

\subsection{\texttt{color\_discrimination\_index}}

Measures how clearly the participant distinguished between the safest colour
(purple) and the most dangerous one (orange), expressed as a normalised
Cohen's~$d$:

\begin{equation}
    d = \frac{\bar{x}_{\text{purple}} - \bar{x}_{\text{orange}}}
             {s_{\text{pooled}}}
    \;,\qquad
    s_{\text{pooled}} = \sqrt{\frac{s^2_{\text{purple}} + s^2_{\text{orange}}}{2}}
\end{equation}

\begin{equation}
    \texttt{color\_discrimination\_index} = \text{clip}\!\left(\frac{d}{2},\; 0,\; 1\right)
\end{equation}

\begin{lstlisting}[style=pythonstyle]
purple_arr = np.array(color_pumps["purple"])
orange_arr = np.array(color_pumps["orange"])

mean_diff  = np.mean(purple_arr) - np.mean(orange_arr)
pooled_std = np.sqrt(
    (np.var(purple_arr, ddof=1) + np.var(orange_arr, ddof=1)) / 2)

if pooled_std == 0:
    return 1.0 if mean_diff > 0 else 0.0

cohens_d = mean_diff / pooled_std
discrimination = np.clip(cohens_d / 2.0, 0.0, 1.0)
\end{lstlisting}

We chose Cohen's~$d$ over a $t$-test because effect sizes are sample-size invariant,
which matters a lot when $N$ per colour can be as low as 2.

Sample session: $0.833$ --- the participant clearly pumped more on purple than
on orange.


\subsection{\texttt{risk\_sensitivity}}

Takes a different angle: Pearson~$r$ between each balloon's capacity
($128$/$32$/$8$) and the participant's actual pump count on that balloon.
This tells us whether the participant scaled their behaviour proportionally to
the available headroom.

\begin{lstlisting}[style=pythonstyle]
capacities, pumps = [], []
for color, ps in color_pumps.items():
    cap = COLOR_PROFILES[color]["max_pumps"]
    for p in ps:
        capacities.append(cap)
        pumps.append(p)

r, _ = stats.pearsonr(capacities, pumps)
\end{lstlisting}

$r > 0.7$ indicates strong sensitivity; $r \approx 0$ a flat strategy; $r < 0$
would be unusual (pumping \emph{more} on dangerous balloons).  Sample session:
$r = 0.678$.

One caveat: a genuinely flat but deliberate strategy also produces $r \approx 0$,
so a low value is ambiguous on its own.  Cross-referencing with
\texttt{color\_discrimination\_index} helps disambiguate.


\subsection{\texttt{risk\_adjustment\_score}}
\label{sec:risk_adj}

Asks a slightly different question: not ``did you differentiate?'' but ``did you
pump the \emph{right amount} for each colour?''  Each colour is scored against
its EV-optimal region, and the three scores are averaged:

\begin{align}
    \text{Purple score}  &= \text{clip}\!\left(\frac{\bar{x}_\text{purple}}{128},\; 0,\; 1\right) \times 100
    &\quad&\text{(more pumps\,$=$\,better)} \\[4pt]
    \text{Orange score}  &= \text{clip}\!\left(1 - \frac{\bar{x}_\text{orange}}{8},\; 0,\; 1\right) \times 100
    &\quad&\text{(fewer pumps\,$=$\,better)} \\[4pt]
    \text{Teal score}    &= \text{clip}\!\left(1 - \frac{|\bar{x}_\text{teal} - 8|}{16},\; 0,\; 1\right) \times 100
    &\quad&\text{(closer to 8\,$=$\,better)}
\end{align}

\begin{lstlisting}[style=pythonstyle]
scores = []
if "purple" in cp and cp["purple"]:
    scores.append(np.clip(np.mean(cp["purple"]) / 128, 0, 1) * 100)
if "orange" in cp and cp["orange"]:
    scores.append(np.clip(1 - np.mean(cp["orange"]) / 8, 0, 1) * 100)
if "teal" in cp and cp["teal"]:
    d = abs(np.mean(cp["teal"]) - 8.0)
    scores.append(np.clip(1 - d / 16, 0, 1) * 100)

risk_adjustment_score = np.mean(scores)
\end{lstlisting}

Range $[0, 100]$.  This is the single best ``risk intelligence'' number because it
penalises both over-pumping on dangerous balloons \emph{and} under-pumping on safe
ones.  It replaced \texttt{explosion\_rate} in the reflection pipeline for exactly
this reason: it reflects what the participant chose, not what the dice rolled.

Sample session: $32.2$ --- below average calibration, driven mainly by low
exploitation of purple's headroom.


% ──────────────────────────────────────────────────────────────────────────
\section{Consistency Metrics}
% ──────────────────────────────────────────────────────────────────────────

\subsection{\texttt{response\_consistency}}

Global coefficient of variation across every intra-pump latency in the session:

\begin{equation}
    \text{CV}_\text{global} = \frac{\sigma(\text{all\_latencies})}{\mu(\text{all\_latencies})}
\end{equation}

Sample session: $1.204$.  On its own this is hard to interpret because it mixes
two very different sources of variability --- motor jitter within a balloon and
strategic tempo shifts between balloons.  The decomposition below is more useful.


\subsection{\texttt{within\_balloon\_consistency}}

For each balloon with $\geq 3$ pumps we compute the CV of that balloon's inter-pump
intervals (filtering out any pauses $> 2$\,s), then average across balloons:

\begin{lstlisting}[style=pythonstyle]
within_cvs = []
for balloon_events in balloons:
    ts = [e.timestamp for e in balloon_events if e.type == "pump"]
    if len(ts) >= 3:
        diffs = np.diff(ts)
        diffs = diffs[diffs < 2000.0]
        if len(diffs) >= 2 and np.mean(diffs) > 0:
            within_cvs.append(np.std(diffs) / np.mean(diffs))

within_balloon_cv = np.mean(within_cvs) if within_cvs else 0.0
\end{lstlisting}

Low values ($< 0.2$) point to steady, rhythmic pumping; high values ($> 0.6$)
suggest hesitation or attentional lapses mid-balloon.


\subsection{\texttt{between\_balloon\_consistency}}

CV of \emph{pump counts} across collected balloons (falling back to all balloons
if fewer than 5 were collected):

\begin{lstlisting}[style=pythonstyle]
collected_counts, all_counts = [], []
for b in balloons:
    p = sum(1 for e in b if e.type == "pump")
    all_counts.append(p)
    terminal = next(
        (e.type for e in reversed(b)
         if e.type in ("collect","explode")), None)
    if terminal != "explode":
        collected_counts.append(p)

arr = np.array(collected_counts if len(collected_counts) >= 5
               else all_counts, dtype=np.float64)
between_cv = np.std(arr) / np.mean(arr)
\end{lstlisting}

\subsubsection*{Putting them together}

\begin{table}[H]
\centering
\begin{tabular}{@{}lll@{}}
\toprule
\textbf{Within CV} & \textbf{Between CV} & \textbf{Interpretation} \\
\midrule
Low  & Low  & Highly consistent --- same strategy, steady execution \\
Low  & High & Strategically variable --- adapts per balloon, but executes cleanly \\
High & High & Erratic --- inconsistent in both strategy and execution \\
High & Low  & Attentionally unstable --- one strategy, but shaky execution \\
\bottomrule
\end{tabular}
\end{table}


% ──────────────────────────────────────────────────────────────────────────
\section{Composite Metrics}
% ──────────────────────────────────────────────────────────────────────────

\subsection{\texttt{adaptive\_strategy\_score}}

A single number meant to summarise ``behavioural intelligence'' in the task.
It combines three equally weighted components --- learning, discrimination, and
calibration --- each contributing up to 33.33 points:

\begin{equation}
\boxed{
\texttt{adaptive\_strategy\_score} =
\underbrace{\frac{\texttt{half\_split\_lr} + 1}{2} \times 33.33}_{\text{Learning}}
\;+\;
\underbrace{\texttt{color\_disc} \times 33.33}_{\text{Discrimination}}
\;+\;
\underbrace{\frac{\texttt{risk\_adj}}{100} \times 33.33}_{\text{Calibration}}
}
\end{equation}

\begin{lstlisting}[style=pythonstyle]
safe = lambda v: 0.0 if np.isnan(v) else v

adaptive_strategy_score = (
    (safe(half_split_lr) + 1) / 2 * 33.33
    + safe(color_discrimination) * 33.33
    + safe(risk_adjustment) / 100 * 33.33
)
adaptive_strategy_score = np.clip(adaptive_strategy_score, 0, 100)
\end{lstlisting}

We gave each component equal weight because there is no strong theoretical reason
to privilege one over the others.  In practice a score above $70$ indicates
someone who learns, differentiates, \emph{and} calibrates well; below~$40$
usually points to either a flat strategy, poor colour discrimination, or
performance that got worse over time.

Sample session: $56.1$ --- decent discrimination pulled the score up, but
weak calibration (especially on purple) held it back.


% ──────────────────────────────────────────────────────────────────────────
\section{Session Validation}
% ──────────────────────────────────────────────────────────────────────────

\subsection{\texttt{session\_valid}}

Before scoring we run five sanity checks:

\begin{enumerate}
    \item \textbf{Minimum balloons.}  Fewer than 15 $\to$ invalid outright;
          15--29 $\to$ warning.
    \item \textbf{Colour balance.}  Each colour should have ${\sim}10$ balloons;
          $< 5$ triggers a warning.
    \item \textbf{Timestamp monotonicity.}  Any $t_i < t_{i-1}$ $\to$ invalid
          (corrupt data).
    \item \textbf{Session speed.}  Under 30\,s for $\geq 15$ balloons is suspiciously
          fast.
    \item \textbf{Pump uniformity.}  If
          $\text{std}(\text{pump\_counts}) < 0.5$ across $\geq 10$ balloons, the
          session likely came from a bot or an auto-clicker.
\end{enumerate}

\begin{lstlisting}[style=pythonstyle,caption={Bot-detection heuristic (Python).}]
pump_counts = [sum(1 for e in b if e.type == "pump")
               for b in balloons]
if len(pump_counts) >= 10 and float(np.std(pump_counts)) < 0.5:
    warnings.append(
        "Suspicious: nearly identical pump counts "
        "across all balloons")
\end{lstlisting}

\subsection{\texttt{session\_warnings}}

A list that accumulates messages from both the validation step and any data-quality
fallbacks that fire during scoring.


% ──────────────────────────────────────────────────────────────────────────
\section{Behavioural Profile (Narrative)}
% ──────────────────────────────────────────────────────────────────────────

After the numbers are computed we also generate a short narrative profile,
primarily for the participant-facing reflection screen.

\subsection{\texttt{risk\_style}}

One of four archetypes, assigned by simple threshold logic:

\begin{table}[H]
\centering
\begin{tabular}{@{}lp{5.5cm}p{5cm}@{}}
\toprule
\textbf{Style} & \textbf{When it fires} & \textbf{Description} \\
\midrule
High-Stakes Risk Taker   & \texttt{rng\_normalized} $> 0.45$ & Pushes limits to the maximum \\
Conservative Safety-Seeker & \texttt{rng\_normalized} $< 0.10$ \textbf{and} \texttt{avg\_adj} $< 10$ & Prioritises safety above all \\
Strategic Risk Taker      & \texttt{risk\_sensitivity} $> 0.7$ & Takes calculated risks only when favourable \\
Balanced Explorer         & default & Healthy mix of caution and exploration \\
\bottomrule
\end{tabular}
\end{table}

Sample session: \emph{High-Stakes Risk Taker}.


\subsection{\texttt{dominant\_traits}}

Additional trait labels, each gated by a single threshold:

\begin{table}[H]
\centering
\begin{tabular}{@{}lp{7.5cm}@{}}
\toprule
\textbf{Trait} & \textbf{Condition} \\
\midrule
Highly Consistent       & within\_CV $< 0.2$ \textbf{and} between\_CV $< 0.4$ \\
Erratic (within-balloon)& within\_CV $> 0.6$ \\
Strategically Variable  & between\_CV $> 1.0$ \\
Adaptive Learner        & \texttt{half\_split\_lr} $> 0.1$ \\
Risk-Averse Learner     & \texttt{half\_split\_lr} $< -0.1$ \\
Impulsive               & \texttt{orange\_avg\_pumps} $> 4.0$ \\
Patient Optimizer       & \texttt{patience\_index} $> 40$ \\
\bottomrule
\end{tabular}
\end{table}

Sample session: \emph{Variable/Erratic}, \emph{Impulsive}.


% ──────────────────────────────────────────────────────────────────────────
\section{Metric Classification Table}
\label{sec:classification}
% ──────────────────────────────────────────────────────────────────────────

\begin{small}
\begin{longtable}{@{}lllc@{}}
\toprule
\textbf{Metric} & \textbf{Type} & \textbf{Data Source} & \textbf{Used in Clustering} \\
\midrule
\endhead
\texttt{average\_pumps\_adjusted}     & Descriptive & Collected only     & No \\
\texttt{avg\_pumps\_all\_balloons}    & Descriptive & All balloons       & No \\
\texttt{explosion\_rate}              & Outcome     & All balloons       & No \\
\texttt{mean\_latency\_between\_pumps}& Temporal    & All balloons       & No \\
\texttt{rng\_normalized\_pumps}       & Behavioural & Collected only     & Yes \\
\texttt{impulsivity\_index}           & Behavioural & Collected orange   & Yes \\
\texttt{patience\_index\_normalized}  & Behavioural & Collected purple   & Yes \\
\texttt{orange\_avg\_pumps}           & Behavioural & Collected orange   & No \\
\texttt{patience\_index}              & Behavioural & Collected purple   & No \\
\texttt{learning\_rate}               & Adaptation  & Collected / colour & No \\
\texttt{half\_split\_learning\_rate}  & Adaptation  & Collected / colour & No \\
\texttt{color\_discrimination\_index} & Calibration & Collected / colour & No \\
\texttt{risk\_sensitivity}            & Calibration & Collected / colour & Yes \\
\texttt{risk\_adjustment\_score}      & Calibration & Collected / colour & No \\
\texttt{response\_consistency}        & Temporal    & All latencies      & No \\
\texttt{within\_balloon\_consistency} & Temporal    & Per-balloon lat.   & No \\
\texttt{between\_balloon\_consistency}& Strategic   & Collected counts   & No \\
\texttt{adaptive\_strategy\_score}    & Composite   & All collected      & No \\
\texttt{session\_valid}               & Validation  & All events         & No \\
\bottomrule
\end{longtable}
\end{small}


% ──────────────────────────────────────────────────────────────────────────
\section{K-Means Clustering Pipeline}
\label{sec:clustering}
% ──────────────────────────────────────────────────────────────────────────

\textbf{Important caveat.}  The results in this section come from a pilot
sample of $N = 12$ participants.  We expect to collect 50--80 participants by the
end of the experiment, at which point the feature set, the number of clusters, and
the profile interpretations will almost certainly change.  What follows is the
current state of the pipeline and its preliminary output --- not a final result.

\subsection{Feature Horizon}

For the pilot run we narrowed the input to five features that seemed to separate
participants most cleanly at this sample size:

\begin{table}[H]
\centering
\begin{tabular}{@{}llp{7cm}@{}}
\toprule
\textbf{Feature} & \textbf{Source} & \textbf{What it captures} \\
\midrule
\texttt{dospert\_financial}           & DOSPERT-30 & Self-reported willingness to take financial risks \\
\texttt{bart\_risk\_sensitivity}      & BART       & Pearson $r$ between balloon capacity and pumping behaviour \\
\texttt{bart\_patience\_normalized}   & BART       & Purple (safe) balloon exploitation, normalised to $[0,1]$ \\
\texttt{bart\_impulsivity\_index}     & BART       & Orange (dangerous) balloon pumping, normalised to $[0,1]$ \\
\texttt{bart\_rng\_normalized\_pumps} & BART       & Capacity-normalised pumping intention across all colours \\
\bottomrule
\end{tabular}
\caption{Features used in the pilot clustering run.  This set will likely expand
once more participants are collected.}
\end{table}


\subsection{Preprocessing}

Every feature is standardised to zero mean and unit variance before clustering:

\begin{lstlisting}[style=pythonstyle,caption={Preprocessing (Python).}]
from sklearn.preprocessing import StandardScaler

# Drop participant_id, keep only clustering features
df_numeric = df_work[CONTINUOUS_COLS].copy()

scaler = StandardScaler()
df_scaled = pd.DataFrame(
    scaler.fit_transform(df_numeric),
    columns=df_numeric.columns,
    index=df_numeric.index,
)
\end{lstlisting}

We also run a quick correlation check and flag any pairs with $|r| > 0.70$
to catch redundancy before it distorts the distance metric.


\subsection{Choosing $k$: Elbow and Silhouette}

We sweep $k \in \{2, 3, 4, 5, 6\}$ and evaluate each solution by inertia
(within-cluster sum of squares) and silhouette score:

\begin{lstlisting}[style=pythonstyle,caption={$k$-selection loop (Python).}]
from sklearn.cluster import KMeans
from sklearn.metrics import silhouette_score

inertias, silhouettes = [], {}
for k in range(2, 7):
    km = KMeans(n_clusters=k, n_init=20, random_state=42,
                max_iter=2000)
    labels = km.fit_predict(X)
    inertias.append(km.inertia_)
    silhouettes[k] = silhouette_score(X, labels)

best_k = max(silhouettes, key=silhouettes.get)
\end{lstlisting}

\begin{figure}[H]
\centering
\includegraphics[width=\textwidth]{01_k_selection.png}
\caption{Elbow plot (left) and silhouette scores (right) for $k = 2 \ldots 6$.
The pilot data favour $k = 2$ with a silhouette of $0.343$.  The elbow curve
does not show a sharp bend, which is expected at $N = 12$ --- there simply
are not enough points to carve out more than two regions cleanly.}
\label{fig:k_selection}
\end{figure}


\subsection{Cluster Profiles}

With $k = 2$ selected, we fit the final model and inspect the $z$-scored centroid
deviations from the overall sample mean:

\begin{lstlisting}[style=pythonstyle,caption={Profile computation (Python).}]
km = KMeans(n_clusters=best_k, n_init=20, random_state=42)
labels = km.fit_predict(df_scaled.values)

centroids = df.groupby("cluster")[features].mean()
z_centroids = (centroids - df[features].mean()) / df[features].std()
\end{lstlisting}

\begin{figure}[H]
\centering
\includegraphics[width=\textwidth]{02_cluster_profiles.png}
\caption{Heatmap of cluster centroids expressed as $z$-scores.  Cluster~0
(``risk-engaged'') sits above the mean on every feature; Cluster~1
(``risk-averse'') mirrors it below.  The strongest separator is
\texttt{bart\_patience\_normalized} ($z = \pm 0.82$), followed by
\texttt{bart\_rng\_normalized\_pumps} ($\pm 0.71$).}
\label{fig:profiles}
\end{figure}

Both clusters contain exactly 6 participants.  The profile is strikingly symmetric
--- Cluster~0 is above the global mean on every feature, Cluster~1 below it ---
which is typical when $k = 2$ on a small, roughly Gaussian sample.  The main
takeaway at this stage is that the split is being driven primarily by BART patience
and overall capacity usage, with DOSPERT financial risk riding along in the same
direction.


\subsection{PCA Projection}

To visualise how well the two clusters separate in reduced space, we project
the standardised features onto their first two principal components:

\begin{lstlisting}[style=pythonstyle,caption={PCA projection (Python).}]
from sklearn.decomposition import PCA

pca = PCA(n_components=2)
X_2d = pca.fit_transform(df_scaled.values)
# PC1: 52.7% variance, PC2: 27.7% (80.4% total)
\end{lstlisting}

\begin{figure}[H]
\centering
\includegraphics[width=0.85\textwidth]{03_pca_projection.png}
\caption{PCA projection of the 12 pilot participants.  PC1 (52.7\% variance)
separates the two clusters cleanly along the left--right axis.  PC2 (27.7\%)
captures remaining within-cluster spread, mostly driven by
\texttt{bart\_impulsivity\_index} (loading $= 0.778$).}
\label{fig:pca}
\end{figure}

The PCA loadings confirm the story from the heatmap:

\begin{table}[H]
\centering
\begin{tabular}{@{}lrr@{}}
\toprule
\textbf{Feature} & \textbf{PC1} & \textbf{PC2} \\
\midrule
\texttt{bart\_patience\_normalized}   & $+0.556$ & $-0.169$ \\
\texttt{bart\_rng\_normalized\_pumps} & $+0.542$ & $+0.196$ \\
\texttt{bart\_risk\_sensitivity}      & $+0.450$ & $-0.527$ \\
\texttt{dospert\_financial}           & $+0.406$ & $+0.225$ \\
\texttt{bart\_impulsivity\_index}     & $+0.172$ & $+0.778$ \\
\bottomrule
\end{tabular}
\caption{PCA loadings.  PC1 is a general ``risk engagement'' axis --- all features
load positively.  PC2 picks up a contrast between impulsivity (positive) and
risk sensitivity (negative), which is interesting: some participants score high
on impulsivity without being particularly risk-\emph{sensitive}.}
\end{table}

Together the first two components explain $80.4\%$ of the total variance, which is
reasonable for five input features.


\subsection{What Changes at Scale}

With 50--80 participants we expect several things to shift:

\begin{itemize}
    \item \textbf{More clusters.}  The $k = 2$ solution is almost certainly a
          side-effect of sample size.  With more data we would hope to see at
          least 3--4 distinct personas (e.g.\ separating ``impulsive risk-taker''
          from ``patient strategist'' from ``cautious across the board'').
    \item \textbf{Wider feature set.}  We plan to reintroduce
          \texttt{bart\_mean\_latency}, \texttt{bart\_between\_consistency}, and
          \texttt{bart\_adaptive\_strategy} once the sample is large enough for
          them to contribute meaningfully.  Additional DOSPERT subscales
          (health/safety, recreational, ethical, social) may also enter.
    \item \textbf{Demographics.}  Binary and nominal demographic variables
          (gender, dorm status, employment, prior task experience) are currently
          excluded from clustering but will be tested as candidate features or
          used for post-hoc profiling.
    \item \textbf{Validation.}  At larger $N$ we can run stability checks
          (bootstrap resampling, split-half consistency) and compare silhouette
          scores more meaningfully.
\end{itemize}


% ──────────────────────────────────────────────────────────────────────────
\section{Appendix: Sample Session Walkthrough}
% ──────────────────────────────────────────────────────────────────────────

Session~\texttt{6fe2f79a-0135-4850-9a5c-f0969be6088c}.

\begin{table}[H]
\centering\small
\begin{tabular}{@{}clccc@{}}
\toprule
\textbf{\#} & \textbf{Colour} & \textbf{Pumps} & \textbf{Outcome} & \textbf{Notes} \\
\midrule
 1 & \textcolor{purpleballoon}{Purple} &  6 & Collected & \\
 2 & \textcolor{tealballoon}{Teal}     &  3 & Exploded  & \\
 3 & \textcolor{tealballoon}{Teal}     &  4 & Exploded  & \\
 4 & \textcolor{tealballoon}{Teal}     &  5 & Collected & \\
 5 & \textcolor{purpleballoon}{Purple} &  5 & Exploded  & \\
 6 & \textcolor{purpleballoon}{Purple} &  3 & Exploded  & \\
 7 & \textcolor{orangeballoon}{Orange} &  1 & Exploded  & first-pump pop \\
 8 & \textcolor{orangeballoon}{Orange} &  4 & Exploded  & \\
 9 & \textcolor{purpleballoon}{Purple} &  6 & Collected & \\
10 & \textcolor{orangeballoon}{Orange} &  5 & Exploded  & \\
11 & \textcolor{tealballoon}{Teal}     &  5 & Collected & \\
12 & \textcolor{purpleballoon}{Purple} & 11 & Collected & \\
13 & \textcolor{tealballoon}{Teal}     &  4 & Exploded  & \\
14 & \textcolor{purpleballoon}{Purple} &  7 & Collected & \\
15 & \textcolor{tealballoon}{Teal}     &  3 & Collected & \\
16 & \textcolor{purpleballoon}{Purple} & 21 & Collected & session high \\
17 & \textcolor{tealballoon}{Teal}     &  2 & Exploded  & \\
18 & \textcolor{orangeballoon}{Orange} &  2 & Exploded  & \\
19 & \textcolor{orangeballoon}{Orange} &  4 & Exploded  & \\
20 & \textcolor{purpleballoon}{Purple} & 15 & Exploded  & \\
21 & \textcolor{tealballoon}{Teal}     &  6 & Exploded  & \\
22 & \textcolor{orangeballoon}{Orange} &  4 & Exploded  & \\
23 & \textcolor{purpleballoon}{Purple} & 16 & Exploded  & \\
24 & \textcolor{orangeballoon}{Orange} &  4 & Exploded  & \\
25 & \textcolor{tealballoon}{Teal}     &  5 & Collected & \\
26 & \textcolor{purpleballoon}{Purple} &  9 & Exploded  & \\
27 & \textcolor{orangeballoon}{Orange} &  2 & Exploded  & \\
28 & \textcolor{orangeballoon}{Orange} &  2 & Exploded  & \\
29 & \textcolor{tealballoon}{Teal}     &  3 & Collected & \\
30 & \textcolor{orangeballoon}{Orange} &  1 & Exploded  & first-pump pop \\
\bottomrule
\end{tabular}
\caption{Complete balloon-by-balloon breakdown.}
\end{table}

\begin{table}[H]
\centering
\begin{tabular}{@{}lr@{}}
\toprule
\textbf{Metric} & \textbf{Value} \\
\midrule
\texttt{total\_balloons}               & 30 \\
\texttt{total\_pumps}                  & 168 \\
\texttt{total\_explosions}             & 20 \\
\texttt{total\_collections}            & 10 \\
\texttt{average\_pumps\_adjusted}      & 7.20 \\
\texttt{explosion\_rate}               & 0.667 \\
\texttt{mean\_latency\_between\_pumps} & 555.4\,ms \\
\texttt{learning\_rate}                & 0.059 \\
\texttt{risk\_adjustment\_score}       & 32.2 \\
\texttt{color\_discrimination\_index}  & 0.833 \\
\texttt{risk\_sensitivity}             & 0.678 \\
\texttt{adaptive\_strategy\_score}     & 56.1 \\
\texttt{response\_consistency}         & 1.204 \\
\texttt{patience\_index}               & 9.9 \\
\texttt{risk\_style}                   & High-Stakes Risk Taker \\
\texttt{dominant\_traits}              & Variable/Erratic, Impulsive \\
\bottomrule
\end{tabular}
\caption{Scored output for the sample session.}
\end{table}


% ──────────────────────────────────────────────────────────────────────────
\section*{References}
% ──────────────────────────────────────────────────────────────────────────

\begin{itemize}

    \item Blais, A.-R.\ \& Weber, E.U.\ (2006).
        A Domain-Specific Risk-Taking (DOSPERT) scale for adult populations.
        \textit{Judgment and Decision Making}, 1(2), 33--47.

    \item Din\c{c}, S.C.\ \& Yava\c{s} Tez, \"O.\ (2019).
        Alana \"ozg\"u risk alma \"ol\c{c}e\u{g}i --- k{\i}sa formu'nun
        (DOSPERT-30) T\"urk\c{c}eye uyarlama \c{c}al{\i}\c{s}mas{\i}
        [The adaptation study into Turkish of the Domain-Specific Risk-Taking
        Scale --- Short Form (DOSPERT-30)].
        \textit{Spor Bilimleri Dergisi}, 30(3), 107--120.
        \url{https://doi.org/10.17644/sbd.471304}

    \item Lejuez, C.W., Read, J.P., Kahler, C.W., et al.\ (2002).
        Evaluation of a behavioral measure of risk taking: The Balloon Analogue
        Risk Task (BART). \textit{Journal of Experimental Psychology: Applied},
        8(2), 75--84.

    \item Weber, E.U., Blais, A.-R., \& Betz, N.E.\ (2002).
        A domain-specific risk-attitude scale: Measuring risk perceptions and
        risk behaviors. \textit{Journal of Behavioral Decision Making}, 15(4),
        263--290.

\end{itemize}


\end{document}